\documentclass[aspectratio=169,12pt]{beamer}
% Maths Lycée — thème éditorial inspiré du portfolio de Yassine.
% Modifier ici les couleurs et la typographie de tous les diaporamas.
\usepackage[T1]{fontenc}
\usepackage[utf8]{inputenc}
\usepackage[french]{babel}
\usepackage{lmodern,amsmath,amssymb,mathtools,tikz,pgfplots,listings,adjustbox}
\usetikzlibrary{arrows.meta,positioning,calc}
\pgfplotsset{compat=1.18,width=9cm,height=4.2cm}
\definecolor{Ink}{HTML}{202124}
\definecolor{Blue}{HTML}{356B9E}
\definecolor{Mist}{HTML}{F6F6F5}
\definecolor{Muted}{HTML}{737579}
\definecolor{Line}{HTML}{E7E7E5}
\setbeamercolor{normal text}{fg=Ink,bg=white}
\setbeamercolor{structure}{fg=Blue}
\setbeamercolor{frametitle}{fg=Ink}
\setbeamercolor{title}{fg=Ink}
\setbeamercolor{block title}{fg=Blue,bg=Mist}
\setbeamercolor{block body}{fg=Ink,bg=Mist}
\setbeamerfont{title}{size=\LARGE,series=\bfseries}
\setbeamerfont{frametitle}{size=\Large,series=\bfseries}
\setbeamersize{text margin left=1cm,text margin right=1cm}
\setbeamertemplate{navigation symbols}{}
\setbeamertemplate{itemize items}[circle]
\setbeamertemplate{enumerate items}[default]
\setlength{\parskip}{7pt}
\newcommand{\slidekind}{COMPRENDRE}
\newcommand{\chapternumber}{01}
\setbeamertemplate{frametitle}{\vspace{0.3cm}{\fontsize{7}{9}\selectfont\color{Blue}\MakeUppercase{\slidekind}}\par\vspace{0.1cm}\insertframetitle\par\vspace{0.08cm}\color{Line}\rule{\textwidth}{0.4pt}\vspace{-0.02cm}}
\setbeamertemplate{footline}{\hspace{1cm}\begin{minipage}{\dimexpr\paperwidth-2cm}\color{Line}\rule{\textwidth}{0.4pt}\par\vspace{0.06cm}{\fontsize{6.5}{8}\selectfont\color{Muted}MATHS LYCÉE\quad / \quad PREMIÈRE SPÉCIALITÉ\hfill\insertframenumber\ /\ \inserttotalframenumber}\end{minipage}\vspace{0.22cm}}
\setbeamertemplate{title page}{
\vspace{0.3cm}
{\fontsize{8}{10}\selectfont\color{Blue}MATHS LYCÉE\quad /\quad 2026--2027}\par
\vspace{0.6cm}
\begin{columns}[c,onlytextwidth]
\begin{column}{.72\textwidth}
{\fontsize{9}{12}\selectfont\color{Muted}CHAPITRE \chapternumber}\par\vspace{0.22cm}
{\usebeamerfont{title}\color{Ink}\inserttitle\par}
\vspace{0.4cm}{\color{Blue}\rule{1.5cm}{2pt}}\par
{\small\color{Muted}Comprendre les notions.\par Justifier chaque étape.}
\end{column}
\begin{column}{.23\textwidth}
\begin{tikzpicture}[scale=.65]
\fill[Mist,rounded corners=6pt](-1.9,-2.1)rectangle(2,2.5);
\draw[Line,step=.5](-1.8,-2)grid(1.9,2.4);
\draw[Muted,->](-1.7,-1.3)--(1.8,-1.3);
\draw[Muted,->](0,-1.6)--(0,2.2);
\draw[Blue,very thick,domain=-1.5:1.5,samples=40]plot(\x,{1.3*\x*\x-1.3});
\fill[Blue](1,0)circle(2.5pt);
\end{tikzpicture}
\end{column}
\end{columns}
\vfill}
\lstdefinestyle{mathspython}{language=Python,basicstyle=\ttfamily\small,keywordstyle=\color{Blue}\bfseries,commentstyle=\color{Muted},stringstyle=\color{Blue},backgroundcolor=\color{Mist},frame=none,xleftmargin=8pt,xrightmargin=8pt,aboveskip=8pt,belowskip=8pt,showstringspaces=false,columns=fullflexible,keepspaces=true,breaklines=true}
\lstnewenvironment{codepython}{\lstset{style=mathspython}}{}
\newcommand{\R}{\mathbb{R}}
\newcommand{\N}{\mathbb{N}}
\newcommand{\E}{\mathbb{E}}
\newcommand{\V}{\mathrm{V}}
\hypersetup{colorlinks=true,urlcolor=Blue}

\title{La fonction exponentielle}
\subtitle{Première spécialité}
\renewcommand{\chapternumber}{07}\author{Maths Lycée}
\date{Année scolaire 2026--2027}
\begin{document}\begin{frame}\titlepage\end{frame}
\begin{frame}{Objectifs du chapitre}\begin{itemize}
\item Définir la fonction exponentielle
\item Utiliser ses propriétés algébriques
\item Résoudre équations et inéquations
\item Dériver une exponentielle affine
\item Relier modèles continus et suites géométriques
\end{itemize}\end{frame}
\renewcommand{\slidekind}{COMPRENDRE}
\begin{frame}[t,fragile]{Une croissance proportionnelle}
Une population fictive augmente d’autant plus vite qu’elle est grande. On cherche une fonction dont la vitesse de variation est proportionnelle à sa valeur.\medskip\par Dans le modèle le plus simple, $f\prime=f$ et la valeur initiale est $f(0)=1$. Cette condition conduit à une nouvelle fonction.
\end{frame}
\renewcommand{\slidekind}{COMPRENDRE}
\begin{frame}[t,fragile]{Définition}
Il existe une unique fonction dérivable sur $\R$ telle que\[f\prime(x)=f(x)\quad\text{et}\quad f(0)=1.\] On l’appelle fonction exponentielle, notée $\exp$. L’existence et l’unicité sont admises ici. On écrit aussi $\exp(x)=e^x$ et $e=\exp(1)\approx2,71828$.
\end{frame}
\renewcommand{\slidekind}{COMPRENDRE}
\begin{frame}[t,fragile]{Une valeur toujours positive}
La fonction exponentielle ne s’annule pas et vérifie $e^x>0$ pour tout réel $x$.\medskip\par Ainsi $e^{-3}$ est positif, même si l’exposant est négatif. Cette propriété permettra de connaître le signe de nombreuses dérivées sans calculer de valeurs approchées.
\end{frame}
\renewcommand{\slidekind}{JUSTIFIER}
\begin{frame}[t,fragile]{Pourquoi l’exponentielle ne s’annule pas}
Poser $h(x)=f(x)f(-x)$, où $f\prime=f$. La règle du produit donne\[h\prime(x)=f(x)f(-x)-f(x)f(-x)=0.\] Donc $h$ est constante et $h(0)=1$. Ainsi $f(x)f(-x)=1$ : aucun facteur ne peut être nul. La continuité et $f(0)=1$ assurent la positivité.
\end{frame}
\renewcommand{\slidekind}{COMPRENDRE}
\begin{frame}[t,fragile]{Additionner les exposants}
Pour tous réels $a,b$,\[e^{a+b}=e^ae^b.\] Une somme dans l’exposant devient un produit de valeurs. Par exemple $e^{x+2}=e^xe^2$. En revanche, $e^{x+2}$ n’est pas égal à $e^x+e^2$.
\end{frame}
\renewcommand{\slidekind}{JUSTIFIER}
\begin{frame}[t,fragile]{Justification de la relation fonctionnelle}
Fixer $a$ et poser $g(x)=e^{a+x}/e^a$. Le dénominateur est non nul. Alors $g\prime=g$ et $g(0)=1$.\medskip\par Par unicité de la fonction de la définition, $g(x)=e^x$. Multiplier par $e^a$ donne $e^{a+x}=e^ae^x$.
\end{frame}
\renewcommand{\slidekind}{COMPRENDRE}
\begin{frame}[t,fragile]{Opposés, différences et puissances}
La relation précédente donne\[e^{-a}=\frac1{e^a},\qquad e^{a-b}=\frac{e^a}{e^b},\qquad(e^a)^n=e^{na}\quad(n\in\mathbb Z).\] Exemple : $e^{3x}e^{-x}/e^2=e^{2x-2}$. Chaque simplification porte sur les exposants du produit ou du quotient.
\end{frame}
\renewcommand{\slidekind}{COMPRENDRE}
\begin{frame}[t,fragile]{Dérivée et variations}
Par définition, $(e^x)\prime=e^x$. Cette dérivée est strictement positive, donc l’exponentielle est strictement croissante sur $\R$.\medskip\par Puisque $e^0=1$, on a $e^x<1$ si $x<0$ et $e^x>1$ si $x>0$. La fonction reste toujours positive.
\end{frame}
\renewcommand{\slidekind}{COMPRENDRE}
\begin{frame}[t,fragile]{La courbe et sa tangente}
\begin{columns}[c,onlytextwidth]\begin{column}{.60\textwidth}\centering\begin{adjustbox}{max width=\linewidth,max height=3.7cm}\begin{tikzpicture}\begin{axis}[width=9cm,height=4.2cm,axis lines=middle,xmin=-2,xmax=1.6,ymin=-.4,ymax=5,xlabel=$x$,ylabel=$y$]\addplot[Blue,thick,domain=-2:1.6,samples=70]{exp(x)};\addplot[gray,domain=-1.4:1.6]{1+x};\addplot[only marks] coordinates {(0,1)};\end{axis}\end{tikzpicture}\end{adjustbox}\end{column}\begin{column}{.36\textwidth} En zéro, la valeur et la pente valent $1$ : la tangente est $y=1+x$.\end{column}\end{columns}
\end{frame}
\renewcommand{\slidekind}{MÉTHODE \& EXEMPLE}
\begin{frame}[t,fragile]{Équations exponentielles}
La stricte croissance donne l’équivalence\[e^u=e^v\quad\Longleftrightarrow\quad u=v.\] Résoudre $e^{2x-1}=e^{x+3}$ revient à résoudre $2x-1=x+3$, donc $x=4$. Cette méthode fonctionne quand on reconnaît deux exponentielles.
\end{frame}
\renewcommand{\slidekind}{MÉTHODE \& EXEMPLE}
\begin{frame}[t,fragile]{Inéquations exponentielles}
Pour tous réels $u,v$,\[e^u\le e^v\quad\Longleftrightarrow\quad u\le v.\] Ainsi $e^{3-2x}\le1$ équivaut à $3-2x\le0$, donc $x\ge3/2$. Le signe s’inverse en divisant par $-2$, pas lors du passage aux exposants.
\end{frame}
\renewcommand{\slidekind}{MÉTHODE \& EXEMPLE}
\begin{frame}[t,fragile]{Une équation ramenée au second degré}
Résoudre $e^{2x}-3e^x+2=0$. Poser $X=e^x>0$.\[X^2-3X+2=(X-1)(X-2)=0.\] On obtient $e^x=1$ ou $e^x=2$. La première donne $x=0$. La seconde se résout numériquement en première : $x\approx0,693$. Le logarithme sera étudié ensuite.
\end{frame}
\renewcommand{\slidekind}{COMPRENDRE}
\begin{frame}[t,fragile]{Dériver une exponentielle affine}
Pour $f(x)=e^{ax+b}$,\[f\prime(x)=a e^{ax+b}.\] Le facteur $a$ ne doit pas être oublié. Pour $f(x)=5e^{-0,2x}$, on obtient $f\prime(x)=-e^{-0,2x}<0$. Le modèle est décroissant.
\end{frame}
\renewcommand{\slidekind}{COMPRENDRE}
\begin{frame}[t,fragile]{Un produit avec une exponentielle}
$f(x)=xe^{-x}$ sur $[0,+\infty[$. Par produit,\[f\prime(x)=e^{-x}-xe^{-x}=(1-x)e^{-x}.\] L’exponentielle est positive : le signe est celui de $1-x$. Le maximum est atteint en $x=1$ et vaut $1/e$.
\end{frame}
\renewcommand{\slidekind}{COMPRENDRE}
\begin{frame}[t,fragile]{Modèle continu et prélèvements réguliers}
Pour $N(t)=N_0e^{kt}$, avec $N_0>0$,\[\frac{N(t+1)}{N(t)}=e^k.\] Aux instants entiers, $u_n=N(n)$ est une suite géométrique de raison $q=e^k$. Le coefficient continu $k$ et le taux discret $q-1$ sont liés mais différents.
\end{frame}
\renewcommand{\slidekind}{COMPRENDRE}
\begin{frame}[t,fragile]{Interpréter un coefficient}
Avec $N(t)=100e^{0,1t}$, on a $N\prime(t)=0,1N(t)$. Le taux instantané est $0,1$ par unité de temps.\medskip\par Sur une unité de temps, le multiplicateur vaut $e^{0,1}\approx1,10517$, soit environ $10,52\,\%$ d’augmentation, et non exactement $10\,\%$.
\end{frame}
\renewcommand{\slidekind}{COMPRENDRE}
\begin{frame}[t,fragile]{Approximation près de zéro}
La tangente fournit $e^h\approx1+h$ pour $h$ proche de zéro. Ainsi $e^{0,02}\approx1,02$.\medskip\par Cette approximation explique pourquoi taux instantané et taux discret sont proches quand le pas est petit. Elle n’autorise pas à remplacer $e^2$ par $3$.
\end{frame}
\renewcommand{\slidekind}{COMPRENDRE}
\begin{frame}[t,fragile]{Une inégalité démontrée}
Poser $g(x)=e^x-x-1$. Alors $g\prime(x)=e^x-1$, négative pour $x<0$ et positive pour $x>0$.\medskip\par Le minimum vaut $g(0)=0$. Donc $e^x\ge1+x$ sur $\R$, avec égalité seulement en zéro. La courbe reste au-dessus de sa tangente en zéro.
\end{frame}
\renewcommand{\slidekind}{À VOUS DE JOUER}\begin{frame}[t]{Vérification des acquis 1}
$e^{a+b}$ est égal à :\begin{enumerate}
\item $e^a+e^b$
\item $e^ae^b$
\item $e^{ab}$
\item $e^a/e^b$
\end{enumerate}\medskip Une seule réponse est correcte. Justifier votre choix.\end{frame}
\renewcommand{\slidekind}{COMPRENDRE LA RÉPONSE}\begin{frame}[t]{Vérification 1 : résolution}
\begin{block}{Réponse 2}$e^ae^b$\end{block}\medskip La somme des exposants correspond au produit des valeurs.\end{frame}
\renewcommand{\slidekind}{À VOUS DE JOUER}\begin{frame}[t]{Vérification des acquis 2}
Près de zéro, $e^h$ est approché par :\begin{enumerate}
\item $h$
\item $1-h$
\item $h^2$
\item $1+h$
\end{enumerate}\medskip Une seule réponse est correcte. Justifier votre choix.\end{frame}
\renewcommand{\slidekind}{COMPRENDRE LA RÉPONSE}\begin{frame}[t]{Vérification 2 : résolution}
\begin{block}{Réponse 4}$1+h$\end{block}\medskip C’est l’expression de la tangente en zéro.\end{frame}
\end{document}
